% ===================================================================
%  TÁBLA Uimh. 4 — An Aibíocht agus an tSeanaois.
%  Traduction 2026 d'apres texto/io, controlee sur texto/fr, texto/en
%  et texto/es. Consigne : voir texto/ga/10-tabla-01.tex.
%
%  LE TABLEAU DE LA PARENTE. L'irlandais compose ses degres par
%  alliance avec « céile » — « athair céile » le beau-pere, « máthair
%  chéile » la belle-mere, « deartháir céile » le beau-frere : le pere,
%  la mere, le frere « de l'epoux ». Un seul procede pour toute la
%  serie, et il se lit encore.
% ===================================================================

%%K t04-apar-1 apar
\VUsaut{4.0mm}
\VUtitre{15.0pt}{60}{TÁBLA Uimh. 4}
\VUsaut{1.6mm}
\VUtitre{11.4pt}{0}{An Aibíocht agus an tSeanaois.}
\VUsaut{3.2mm}

%%K t04-tit-1 sub
\VUcentre{11.4pt}{0}{\textit{An Chéad Radharc.}}
\VUsaut{1.6mm}
\VUtitre{11.4pt}{0}{Féasta na bainise.}
\VUsaut{2.6mm}

%%K t04-tit-2 sub
\VUpk{10.2pt}{40}{An Fómhar.}
\VUsaut{3.0mm}

%%K t04-c1-01-1 p
1. --- Tá na daoine óga tar éis fás. Tá duine acu, Eoin, ag pósadh. Táimid ag
\VUgras{fhéasta na bainise}, a bhailigh ag an mbord céanna teaghlaigh an fhir
nuaphósta agus na mná nuaphósta.

%%K t04-c1-02-1 p
2. --- Tá an \VUgras{Fómhar} ann, mí \VUgras{Mheán Fómhair}. Níor bhain gaoth
fhuar \VUgras{Dheireadh Fómhair} ná na \VUgras{Samhna} an duilliúr glas de na
páirceanna fós. Tá aer an tráthnóna bog agus cumhra; sin an fáth a dtugtar an
féasta faoin \VUgras{vearanda} \textsuperscript{(8)}, a bhreathnaíonn ar an
ngairdín.

%%K t04-c1-03-1 p
3. --- I bhfad uainn, ar an \VUgras{gcnoc} \textsuperscript{(3)}, tá teach
thuismitheoirí Eoin, \VUgras{teach mór} galánta \textsuperscript{(1)}, folaithe sa
ghlasra; clúdaíonn fíonghoirt áille, a thugann fíon den scoth, fána an chnoic. Tugann
an fíonadóir aire dóibh agus déanann sé peataireacht orthu.

%%K t04-c1-04-1 p
4. --- Os cionn na bhfíniúna eitlíonn ealta \VUgras{patraiscí}
\textsuperscript{(2)}, scanraithe ag teacht an \VUgras{mhaoir choille}
\textsuperscript{(5)}. Agus siúd é, lena \VUgras{ghunna} faoina ascaill agus lena
\VUgras{mhála seilge} ar a ghualainn, ag cuardach na \VUgras{dtor}
\textsuperscript{(4)} lena \VUgras{mhadra} \textsuperscript{(6)}. Coinníonn sé
faire ar an eastát agus ní ligeann sé do na póitseálaithe an \VUgras{fhiabheatha} a
scriosadh.

%%K t04-c1-05-1 p
5. --- Ón taobh amuigh casann \VUgras{fíniúin scáileáin} timpeall na
\VUgras{vearanda}, agus crochtar uaithi \VUgras{triopaill} throma
\textsuperscript{(7)}.

%%K t04-c1-06-1 p
6. --- Is iad na bláthanna áille fómhair a mhaisíonn an \VUgras{bord}
\textsuperscript{(16)} ná na \VUgras{criosaintéimí} \textsuperscript{(9)}; chuir
\VUgras{athair} na mná nuaphósta \textsuperscript{(10)} ceann acu i lúbóg a chóta.

%%K t04-c1-07-1 p
7. --- Tá an féasta ag druidim chun deiridh. Tá an \VUgras{seirbhíseach}
\textsuperscript{(11)} ag tosú ar mhiasa an mhilseoige a thabhairt isteach agus
baileoidh sé gan mhoill na codanna den ghléas boird nach bhfuil ag teastáil a
thuilleadh --- na \VUgras{spúnóga} \textsuperscript{(14)}, na \VUgras{foirc}
\textsuperscript{(12)}, na \VUgras{sceana} \textsuperscript{(13)}, na
\VUgras{miasa} \textsuperscript{(15)}.

%%K t04-tit-3 sub
\VUpk{10.2pt}{40}{An Teaghlach.}
\VUsaut{3.0mm}

%%K t04-c1-08-1 p
8. --- Tá cuma áthasach ar chách, agus taispeánann an gáire lena mbreathnaíonn
\VUgras{máthair} an fhir nuaphósta \textsuperscript{(17)} ar a clann cé chomh sona
atá sí.

%%K t04-c1-09-1 p
9. --- Nascann an pósadh seo dhá theaghlach shaibhre sa cheantar. Is é Eoin, an
\VUgras{fear céile} \textsuperscript{(18)}, \VUgras{mac} úinéara mhóir talún, agus
tá sé ag éirí ina \VUgras{chliamhain} ag monaróir rathúil.

%%K t04-c1-10-1 p
10. --- Tá na \VUgras{céilí} óg, agus mar sin féin bhí siad
\VUgras{geallta} le fada. Tá an \VUgras{bhean chéile} óg \textsuperscript{(19)},
Marta, gleoite ina \VUgras{gúna bán}, maisithe le bláth na n-oráistí.

%%K t04-c1-11-1 p
11. --- Tá a \VUgras{hathair céile} \textsuperscript{(20)} sásta le
\VUgras{bean mhic} chomh maith sin, agus déanann \VUgras{máthair chéile} Eoin
\textsuperscript{(21)} gáire agus í ag féachaint ar a \VUgras{hiníon} chomh
hálainn sin.

%%K t04-c1-12-1 p
12. --- Ag ceann an bhoird suíonn na \VUgras{haíonna} eile \textsuperscript{(22)}:
\VUgras{gaolta, cairde, col ceathracha}. Freastalaíonn an \VUgras{príomhfhreastalaí}
\textsuperscript{(23)} orthu go tapa, agus naipcín ar a lámh.

%%K t04-tit-4 sub
\VUpk{10.2pt}{40}{Na Cairde.}
\VUsaut{3.0mm}

%%K t04-c1-13-1 p
13. --- Ar dheis an bhoird tá \VUgras{iníon uasal} \textsuperscript{(24)},
\VUgras{cara} na mná nuaphósta, ansin deartháir na mná nuaphósta, a
\VUgras{fear finné} \textsuperscript{(25)}. Bíonn sé ag comhrá lena
\VUgras{chailín coimhdeachta} \textsuperscript{(26)}, deirfiúr an fhir nuaphósta,
atá ag éirí ina \VUgras{deirfiúr chéile} ag Marta tríd an bpósadh seo. Is bean óg
chairdiúil í. Tá \VUgras{seoda} áille aici, \VUgras{muince péarlaí} agus
\VUgras{bráisléad} óir.

%%K t04-c1-14-1 p
14. --- In aice leo tá an duine uasal a d'éirigh ina sheasamh agus a d'ardaigh a
ghloine --- is é \VUgras{cara} an teaghlaigh \textsuperscript{(27)} é. Is duine de
\VUgras{fhinnéithe an fhir nuaphósta} é. Deir sé \VUgras{sláinte}, ólann sé ar son
\VUgras{sláinte} na lánúine óige agus guíonn sé sonas fada orthu. Tá sé i gculaith
ghalánta: casóg fhada agus \VUgras{bróga snasta} \textsuperscript{(28)}. Tugann sé
a lámh don \VUgras{tseanmháthair} \textsuperscript{(29)}, atá ina
\VUgras{baintreach}.

%%K t04-c1-15-1 p
15. --- Is é an fear óg sa \VUgras{chasóg fhada} \textsuperscript{(31)}, le
croiméal caol dubh, \VUgras{deartháir céile} Eoin \textsuperscript{(30)}. Is
innealtóir mór le rá é. Suíonn sé in aice le \VUgras{bean uasal} an-ghalánta
\textsuperscript{(33)}, \VUgras{deirfiúr is sine} Mharta agus máthair an chailín
ghleoite atá ag teacht ar láimh lena col ceathrar chun bláth a thabhairt agus chun
\VUgras{tráthnóna maith} a \VUgras{ghuí}. Tá \VUgras{cluasáin} áille agus
\VUgras{gruaig chóirithe} ghalánta ag an mbean uasal seo.

%%K t04-c1-16-1 p
16. --- Tá an col ceathrar beag, chomh greannmhar sin ina \VUgras{bhlús}
\textsuperscript{(36)} agus ina \VUgras{bhrístí} leathana \textsuperscript{(37)},
an-truamhéalach. Is \VUgras{dílleachta} é \textsuperscript{(35)}. Chaill sé a
athair agus a mháthair agus é an-óg. Is é a uncail a \VUgras{chaomhnóir}
\textsuperscript{(38)}, agus d'fhan sé gan phósadh chun an leanbh bocht a thógáil.
Is duine maith é. Tá sé sách maol agus an-ghearr-radharcach; caitheann sé
\VUgras{pionsnaí}.

%%K t04-tit-5 sub
\VUsaut{2.4mm}
\VUpk{10.2pt}{40}{An Milseog.}
\VUsaut{3.0mm}

%%K t04-c1-17-1 p
17. --- Ar chúl na n-aíonna, ar bhord clúdaithe le \VUgras{éadach boird}, tá an
\VUgras{mhilseog} \textsuperscript{(39)} leagtha amach. I vása mór tá torthaí:
\VUgras{péitseoga} \textsuperscript{(40)}, \VUgras{piorraí}
\textsuperscript{(41)}, \VUgras{úlla} \textsuperscript{(42)}, \VUgras{haibreoga}
\textsuperscript{(43)} agus \VUgras{fíonchaora} \textsuperscript{(44)}. Lena thaobh
--- \VUgras{anainn} \textsuperscript{(45)}, \VUgras{císte} álainn
\textsuperscript{(46)}, \VUgras{buidéil licéir} \textsuperscript{(48)} agus
\VUgras{seaimpéin} \textsuperscript{(49)}, agus cruacha \VUgras{plátaí} glana
\textsuperscript{(47)}.

%%K t04-c1-18-1 p
18. --- Osclaíonn an \VUgras{buidéalaí} \textsuperscript{(50)} \VUgras{buidéal}
\textsuperscript{(52)} seanfhíona lena \VUgras{chorcscriú} \textsuperscript{(51)}.
Tá an \VUgras{corc} \textsuperscript{(53)} an-docht, is cosúil. Tá
\VUgras{naipcín} \textsuperscript{(54)} ar lámh an bhuidéalaí seo.

%%K t04-c1-19-1 p
19. --- Tar éis an fhéasta rachaidh na fir chuig an seomra tobac, agus rachaidh na
mná chun iad féin a ghléasadh don rince.

%%K t04-tit-6 sub
\VUsaut{5.0mm}
\VUcentre{11.4pt}{0}{\textit{An Dara Radharc.}}
\VUsaut{1.6mm}
\VUpk{11.4pt}{40}{Lá ainm an tseanathar.}
\VUsaut{2.6mm}

%%K t04-tit-7 sub
\VUpk{10.2pt}{40}{An Chuairt.}
\VUsaut{3.0mm}

%%K t04-c2-01-1 p
1. --- Tá deich mbliana caite. Tá tuismitheoirí Eoin tar éis dul in aois agus is
mór acu a dtriúr \VUgras{gharpháistí}: beirt \VUgras{gharmhac}
\textsuperscript{(2)} agus \VUgras{ghariníon} \textsuperscript{(3)}. Ó tharla gurb é
inniu \VUgras{lá ainm an tseanathar}, tá na páistí tagtha chun comhghairdeas a
dhéanamh leis.

%%K t04-c2-02-1 p
2. --- Tá Peadar beag an-óg fós; níl sé ach trí bliana d'aois. Feiceann sé an
\VUgras{cat} \textsuperscript{(1)} ag teacht chuige. Ba mhaith leis a
\VUgras{fhionnadh} álainn síodúil agus a \VUgras{eireaball} fada a chuimilt; ní
thagann sé isteach ina cheann go bhfuil faoi na \VUgras{lapaí} galánta sin
crúba géara olca, a d'fhéadfadh scríob a dhéanamh air.

%%K t04-c2-03-1 p
3. --- Tá a dheirfiúr Gabrielle, atá á choinneáil ar láimh, i bhfad níos dáiríre,
agus tá Máirtín, ina \VUgras{chóta mór} \textsuperscript{(4)} agus ina
\VUgras{chrios} leathair \textsuperscript{(5)}, ar nós duine fásta, cé go bhfuil
\VUgras{stocaí} \textsuperscript{(6)} le feiceáil faoina bhrístí gearra.

%%K t04-c2-04-1 p
4. --- Chuir a n-athair air a \VUgras{chóta mór} tiubh geimhridh
\textsuperscript{(7)} le \VUgras{bóna fionnaidh} \textsuperscript{(8)}, agus tá
scaif ina \VUgras{phóca} \textsuperscript{(9)}. Ar a bhean tá hata feilte le
\VUgras{cleití} \textsuperscript{(10)}, \VUgras{seaicéad} \textsuperscript{(11)},
\VUgras{bó} fionnaidh \textsuperscript{(12)} agus \VUgras{muif}
\textsuperscript{(13)}.

%%K t04-c2-05-1 p
5. --- Tá sé an-fhuar, mar tá an geimhreadh tagtha. Bhí an \VUgras{sneachta}
\textsuperscript{(19)} ag titim ar feadh an lae, agus tá díonta na dtithe go
hiomlán bán. Leis an \VUgras{oíche} téann an sioc i neart. Sa spéir soilsíonn an
\VUgras{ghealach} \textsuperscript{(17)} agus na \VUgras{réaltaí}
\textsuperscript{(18)}, ag polladh an \VUgras{dorchadais}.

%%K t04-tit-8 sub
\VUsaut{4.2mm}
\VUpk{10.2pt}{40}{An Bhreoiteacht.}
\VUsaut{3.0mm}

%%K t04-c2-06-1 p
6. --- Tá an \VUgras{dall} bocht \textsuperscript{(20)}, atá ag trasnú na cearnóige
os comhair na \VUgras{cógaslainne} \textsuperscript{(16)}, á threorú ag a
\VUgras{phúdal} \textsuperscript{(21)}, an-mhí-ámharach. Tá Iustín, an
\VUgras{cailín aimsire} \textsuperscript{(14)}, ag tabhairt ón gcistin
\VUgras{cupán} \textsuperscript{(15)} tae luibhe an-te, milsithe go flúirseach.

%%K t04-c2-07-1 p
7. --- Tá an \VUgras{tseanmháthair} \textsuperscript{(23)}, ar mhaith léi a
\VUgras{spéaclaí} \textsuperscript{(26)} a thógáil ón mbord chun an
\VUgras{oideas} \textsuperscript{(27)} a léamh, atá díreach scríofa ag an
\VUgras{dochtúir} \textsuperscript{(25)}, ag éirí as an gcathaoir uilleann, ag
brath ar a \VUgras{bata} \textsuperscript{(24)}.

%%K t04-c2-08-1 p
8. --- Fulaingíonn sí go minic \VUgras{tinnis} beaga, \VUgras{meadhráin},
\VUgras{slaghdáin} agus \VUgras{tinneas fiacaile}; tá a cosa lag, agus níl a súile
go maith a thuilleadh. Agus mar sin féin tá sí sásta a bheith ag magadh faoina
sean-\VUgras{dhochtúir}, atá ag iarraidh \VUgras{cógas} \textsuperscript{(28)} a
ordú di i gcónaí. Inniu tá sí an-sona go bhfuil na gaolta óga timpeall uirthi.

%%K t04-c2-09-1 p
9. --- Sa seomra atá dúnta go maith tá sé cluthar. Sa \VUgras{teallach} marmair
\textsuperscript{(29)} tá \VUgras{tine mhaith} ag lasadh \textsuperscript{(30)},
agus tá sé go deas sa \VUgras{solas} bog ón \VUgras{lampa}
\textsuperscript{(31)} atá díreach lasta.

%%K t04-tit-9 sub
\VUsaut{4.2mm}
\VUpk{10.2pt}{40}{An Seanathair.}
\VUsaut{3.0mm}

%%K t04-c2-10-1 p
10. --- Rinne fiú an \VUgras{seanathair} \textsuperscript{(33)} dearmad go raibh sé
tinn, go gcoinníonn an \VUgras{daitheacha} ceangailte den \VUgras{chathaoir
uilleann} \textsuperscript{(34)} é agus go raibh \VUgras{fiabhras agus laige} air
inné féin. Ní cuimhin leis a thuilleadh go raibh \VUgras{niúmóine} air an geimhreadh
seo caite agus gur chuir \VUgras{casacht} chrua phianmhar go mór air.

%%K t04-c2-10-2 p
Agus mar sin féin is ar éigean a tháinig sé chuige féin. Anois tá sé níos fearr.
Ach tá an chos, atá aige ar \VUgras{philiúr} \textsuperscript{(38)} curtha ar
\VUgras{stólín} íseal \textsuperscript{(39)}, ag cur pian air fós.

%%K t04-c2-11-1 p
11. --- Nuair a bhíonn sé fillte go cúramach ina \VUgras{fhallaing} the
\textsuperscript{(35)} le \VUgras{cnaipí} móra móruisce \textsuperscript{(36)},
agus nuair a bhíonn ina aice, chun an nuachtán a léamh dó, an \VUgras{dílleachta}
bheag \textsuperscript{(40)} lena \VUgras{caidhp} bhán, lena \VUgras{gúna} gorm
\textsuperscript{(42)} agus lena \VUgras{cába} \textsuperscript{(41)}, ní
ghearánann sé.

%%K t04-c2-12-1 p
12. --- Gach bliain, nuair a thugann an \VUgras{féilire} \textsuperscript{(43)} ar
ais an chéad \VUgras{lá} de mhí na \VUgras{Nollag}, fanann sé le mífhoighne lena
\VUgras{gharpháistí}, mar tá a fhios aige nach ndéanfaidh siad dearmad ar an
\VUgras{lá} tábhachtach seo.

%%K t04-c2-13-1 p
13. --- Cuimhníonn sé le mothúchán ar na laethanta i bhfad ó shin nuair a théadh sé
féin chun comhghairdeas a dhéanamh leis an \VUgras{marascal} cáiliúil
impiriúil, a bhfuil a \VUgras{phortráid} \textsuperscript{(44)} ar crochadh ar an
mballa agus atá ina \VUgras{shin-seanathair} ag na páistí.
\VUsaut{6.0mm}
\VUfilet{22mm}
